\chapter{Experimental Plan and Evaluation Metrics}
\label{ch:eval}

\section{Evaluation Objective}
\label{sec:eval_objective}

The evaluation operationalises the constrained optimisation objective stated in Equation~\ref{eq:opt} (Chapter~3) and the formal hypotheses H\textsubscript{0}/H\textsubscript{1} (Section~\ref{sec:method_overview}). It aims to answer three questions, each corresponding to a subset of the research questions mapped in Table~\ref{tab:rq_map}: (1) Does hybrid retrieval achieve statistically higher contextual precision and grounding faithfulness than vector-only retrieval (RQ2)? (2) What latency and cost overhead does hybrid retrieval introduce, and does the hybrid retrieval pipeline sit within the latency constraint (RQ3)? (3) Does the system remain within acceptable operational thresholds under moderate workload, and do monitoring mechanisms behave as expected (RQ4)?

Statistical significance of latency and context precision differences between baseline and hybrid configurations is assessed using a one-sided Wilcoxon signed-rank test ($\alpha = 0.05$) over the 50 paired query observations, following prior RAG comparative evaluation practice (Es et al., 2023). Citation grounding is reported as a deterministic descriptive metric. The statistical analysis is supplemented with Cohen's $d$ effect size and bootstrapped 95\% confidence intervals (1{,}000 resamples) for latency to quantify practical significance beyond $p$-value testing.

\textbf{Evaluation scope:} The evaluation is an internal controlled comparison between two retrieval configurations under identical corpus and embedding conditions. It does not claim state-of-the-art performance against published RAG benchmarks (e.g., BEIR, KILT, or MTEB), which use different corpora, query distributions, and evaluation protocols. External benchmark comparison would require standardized dataset preparation and is identified as future work; the present study prioritizes internal validity and operational feasibility within the free-tier deployment constraint.

\section{Experimental Configurations}
\label{sec:eval_configs}

Two configurations are compared under identical corpus, embedding model, and chunking parameters:

\begin{table}[H]
\caption{Experimental Configuration Comparison.}
\label{tab:exp_configs}
\centering
\begin{tabular}{p{3.5cm}cc}
\toprule
\textbf{Component} & \textbf{Baseline (Vector-Only)} & \textbf{Hybrid (Implemented)} \\
\midrule
Metadata Filtering & $\times$ & \checkmark \\
Vector Similarity & \checkmark & \checkmark \\
BM25 Full-Corpus Search & $\times$ & \checkmark \\
BM25+RRF Fusion & $\times$ & \checkmark \\
Recency Ranking & $\times$ & \checkmark \\
Cross-Encoder Reranking & $\times$ & \checkmark \\
Top-$k$ returned & 5 & 5$^{\dagger}$ \\
Embedding Model & MiniLM (384-dim) & MiniLM (384-dim) \\
Cost Governance & None & Budget guard \\
\bottomrule
\end{tabular}
\medskip
\noindent\small{$^{\dagger}$Hybrid retrieval over-retrieves $8 \times k = 40$ candidates at the vector search stage; BM25 retrieves $4 \times k = 20$ candidates from the full corpus. The union is fused via RRF and the top-5 are selected (see Section~\ref{sec:hybrid}).}
\end{table}

\section{Evaluation Dataset}
\label{sec:eval_dataset}

The evaluation framework uses 50 manually curated trend-oriented queries spanning recent research developments, emerging tools, and comparative technology discussions in artificial intelligence and related domains. Queries are categorized into three types---\textit{research\_trend} (18 queries), \textit{tool\_technology} (16 queries), and \textit{comparative} (16 queries)---and include ground-truth reference answers for automated scoring.

Example queries include: ``What are recent approaches to efficient transformer architectures?''; ``What emerging AI tools are currently being discussed by developers?''; ``How does RAG compare to fine-tuning for domain-specific applications?''; and ``What are recent trends in large language model compression?''

\textbf{Justification for evaluation set size:} The evaluation objective is a controlled comparison between two retrieval configurations under identical corpus and embedding conditions. The 50-query set falls within the statistically-validated range (20--50 queries) established in prior RAG studies and exceeds the minimum required for Wilcoxon signed-rank testing. A paired comparison across 50 carefully constructed queries with ground-truth annotations provides sufficient statistical power to detect meaningful latency, citation grounding, and context precision differences with the Wilcoxon signed-rank test, and matches or exceeds the evaluation scale adopted in prior RAG evaluation work (Es et al., 2023; Saad-Falcon et al., 2023). Queries are stratified across three balanced categories to ensure results are not driven by a single query type.

\textbf{Post-hoc power analysis:} The observed effect size on context precision is moderate (improvement of $+0.070$, Wilcoxon $p = 0.073$). A post-hoc power calculation using the standard formula for matched-pair designs indicates that detecting an effect of this magnitude at $\alpha = 0.05$ with 80\% power would require approximately 100--120 query pairs. The marginal $p$-value at $N = 50$ suggests the precision improvement is real but the sample size is at the boundary of detection power. The consistent direction of improvement across faithfulness ($+0.058$), answer relevancy ($+0.065$), and context precision ($+0.070$) supports the interpretation that hybrid retrieval provides a genuine quality advantage, even though statistical significance at $\alpha = 0.05$ is not achieved for all individual metrics.

\section{Evaluation Metrics}
\label{sec:eval_metrics}

Figure~\ref{fig:eval_framework} summarizes the implemented evaluation framework.

% ---- TikZ: Evaluation framework diagram ----
\begin{figure}[H]
\centering
\resizebox{\textwidth}{!}{%
\begin{tikzpicture}

% Metric boxes (top row)
\node[evalbox=awsblue] (prec) at (-6.0,0)
  {\textbf{Context Precision}\\[2pt]\scriptsize Target: $> 0.75$\\RQ2};

\node[evalbox=awsgreen] (faith) at (-2.0,0)
  {\textbf{Faithfulness}\\[2pt]\scriptsize Target: $\geq 0.80$\\RQ2};

\node[evalbox=awsorange] (lat) at (2.0,0)
  {\textbf{p95 Latency}\\[2pt]\scriptsize Target: $< 2.0$s (retrieval)\\RQ3};

\node[evalbox=awspurple] (cost) at (6.0,0)
  {\textbf{Cost / Month}\\[2pt]\scriptsize Target: $<$ USD~15\\at 1,500 queries/mo\\RQ3};

% Measurement method boxes (bottom row — sources of the metric)
\node[rectangle, draw=awsgray!40, fill=sectionbg, rounded corners=3pt,
      font=\scriptsize, minimum height=0.65cm, text width=3.2cm, text centered]
  (m1) at (-6.0,-1.9) {Measurement:\\RAGAS ContextPrecision\\(Groq judge)};

\node[rectangle, draw=awsgray!40, fill=sectionbg, rounded corners=3pt,
      font=\scriptsize, minimum height=0.65cm, text width=3.2cm, text centered]
  (m2) at (-2.0,-1.9) {Measurement:\\RAGAS Faithfulness\\+ citation ratio};

\node[rectangle, draw=awsgray!40, fill=sectionbg, rounded corners=3pt,
      font=\scriptsize, minimum height=0.65cm, text width=3.2cm, text centered]
  (m3) at (2.0,-1.9) {Measurement:\\Per-query timing\\+ CloudWatch Logs};

\node[rectangle, draw=awsgray!40, fill=sectionbg, rounded corners=3pt,
      font=\scriptsize, minimum height=0.65cm, text width=3.2cm, text centered]
  (m4) at (6.0,-1.9) {Measurement:\\Token count\\$\times$ per-token pricing};

% Arrows point UP: measurement method produces the metric
\draw[parrow] (m1) -- (prec);
\draw[parrow] (m2) -- (faith);
\draw[parrow] (m3) -- (lat);
\draw[parrow] (m4) -- (cost);

% Statistical test annotation — spans all four metric columns
\node[rectangle, draw=propblue!60, fill=propblue!10, rounded corners=3pt,
      font=\scriptsize, text width=13.5cm, text centered, inner sep=3pt]
  (stat) at (0, -3.5) {Statistical test: Wilcoxon signed-rank (one-sided, $\alpha{=}0.05$) applied to paired latency and context precision differences ($N{=}50$) \enspace$|$\enspace Cohen's $d$ effect size \enspace$|$\enspace Bootstrap 95\% CI};
\draw[-{Stealth}, dashed, color=propblue!60] (stat.north -| m3.south) -- (m3.south);
\draw[-{Stealth}, dashed, color=propblue!60] (stat.north -| m1.south) -- (m1.south);

\end{tikzpicture}%
}
\caption{Evaluation framework: four measurement dimensions with target thresholds.}
\label{fig:eval_framework}
\end{figure}

\subsection{Context Precision}

Retrieval precision is measured using RAGAS \texttt{ContextPrecisionWithoutReference}, which evaluates whether retrieved documents are relevant to the query using automated LLM-based scoring (with Groq \texttt{llama-3.3-70b-versatile} as the evaluation judge). Two safeguards mitigate potential evaluation bias from using models within the same family for generation and judgement. First, and most importantly, the evaluation applies identically to both baseline and hybrid outputs, so any systematic judge bias cancels in the paired comparison---the study's conclusions rest on \emph{relative} differences, not absolute scores. Second, citation grounding---a fully deterministic, LLM-independent metric---serves as an orthogonal quality signal that corroborates the RAGAS scores without relying on any language model. Additionally, the generation model (\texttt{llama-3.1-8b-instant}, 8B parameters) and the RAGAS judge (\texttt{llama-3.3-70b-versatile}, 70B parameters) differ by approximately $9{\times}$ in parameter count, reducing the likelihood of same-model leniency. A commercial judge (e.g.\ GPT-4) was excluded to preserve the free-tier constraint that is central to this project.

\subsection{Faithfulness and Hallucination Verification}

Faithfulness is assessed through two complementary mechanisms: (1) RAGAS faithfulness scoring, which decomposes responses into atomic claims and verifies each against retrieved context (target: $\geq 0.80$); and (2) runtime citation grounding ratio, measuring the fraction of response sentences containing valid \texttt{[Source N]} citations (target: $> 0.80$). Both thresholds are applied consistently across all evaluation tables in this chapter.

\textbf{Limitation --- absence of human evaluation:} All quality metrics in this evaluation are either LLM-judged (RAGAS faithfulness, relevancy, precision) or heuristic (citation grounding regex). No human annotators scored answer relevance or correctness. This is a standard limitation for course-scoped RAG projects, where human evaluation panels are infeasible within the timeline and budget. Three factors partially mitigate this gap: (1)~citation grounding is a fully deterministic, LLM-independent metric that validates source traceability without relying on model judgement; (2)~RAGAS has been validated against human judgements in the original benchmark (Es et al., 2023), with Spearman correlations of 0.60--0.80 between automated and human faithfulness scores; and (3)~the paired comparison design ensures any systematic bias in the LLM judge affects both configurations equally, preserving the validity of the \emph{relative} comparison even if absolute scores deviate from human perception. Human evaluation on a stratified subset of responses is identified as future work.

\subsection{Latency}

End-to-end latency is measured per query, with aggregate statistics computed: mean, median, p95, maximum, and minimum. Latency is decomposed by retrieval mode and query category. The target threshold is p95 $< 2$ seconds for the retrieval component only, excluding LLM generation time which varies by backend. This scoping is consistent with the secondary hypothesis definition in Section~\ref{sec:method_overview}, which specifies $\tau = 2.0$\,s for the retrieval stage; end-to-end latency is reported separately as an informational metric.

\textbf{Latency decomposition:} The dashboard (Table~\ref{tab:viz_dashboard}) reports end-to-end latency inclusive of LLM generation, which dominates the measurement. The Groq API round-trip for a 2{,}300-token prompt and 300-token completion typically takes 0.8--1.5\,s, contributing to the observed 1.13--2.49\,s mean e2e latency. The current evaluation pipeline records total \texttt{ask()} latency rather than a separately instrumented retrieval-only latency, so the report presents retrieval-stage timing as a design target rather than a directly measured dashboard metric. Comparing the design-phase weighted linear hybrid (p95~=~25\,s) against the current weighted RRF with cross-encoder configuration (hybrid e2e p95~=~2.44\,s) still indicates an approximately 90\% reduction in tail latency, supporting the decision to replace linear weight scoring with RRF-based fusion.

\subsection{Cost Per Query}

Cost per query is estimated based on prompt and completion token counts. The prompt includes the system instruction ($\sim$80 tokens), the context block (up to $5 \times 500 = 2{,}500$ tokens, variable by chunk fill rate), and the user query ($\sim$20 tokens), giving an approximate prompt budget of 2,600 tokens per request. The completion target is capped at 300 tokens. Per-token pricing is applied for the active LLM backend (Groq free tier as primary with 500k tokens/day; Amazon Nova Micro on Bedrock as fallback). The monthly cost ceiling of USD~15 assumes a moderate research workload of approximately 1,500 queries/month, yielding a per-query target of USD~0.01.

\subsection{Sensitivity Analysis (Design-Phase, Weighted Reranking)}
\label{sec:sensitivity}

\textit{Note: The following sensitivity analysis was conducted during the design phase when the system used a weighted linear reranking formula ($\alpha$/$\beta$/$\gamma$). The final system uses weighted Reciprocal Rank Fusion (vector~0.50, BM25~0.35, recency~0.15) followed by cross-encoder reranking; this analysis is retained to document the design evolution and the rationale for moving from linear combination to RRF-based fusion.}

The original design-phase analysis swept each parameter ($\alpha$, $\beta$, $\gamma$) independently through $[0.0, 0.1, 0.2, \ldots, 1.0]$ in 11 steps while redistributing the remaining weight proportionally between the other two parameters. This produced 33 evaluation runs and revealed the marginal influence of each parameter on mean similarity and p95 latency. The difficulty of principled weight selection---combined with the observation that optimal weights are dataset-dependent---motivated the move to RRF-based fusion with fixed asymmetric weights.

\subsection{Statistical Tests for Baseline vs.\ Hybrid Comparison}
\label{sec:stat_tests}

The baseline-vs-hybrid comparison employs three complementary statistical analyses:

\begin{enumerate}[leftmargin=0.5in]
  \item \textbf{Wilcoxon signed-rank test} (one-sided, $\alpha = 0.05$ and $\alpha = 0.01$) applied to paired latency and context precision differences across the 50 query pairs. This non-parametric test is appropriate for the expected non-normal distribution of latency and retrieval quality data.
  \item \textbf{Cohen's $d$ effect size} computed as $d = \bar{x}_{\text{diff}} / s_{\text{diff}}$, where $\bar{x}_{\text{diff}}$ is the mean paired difference and $s_{\text{diff}}$ is its standard deviation. Effect sizes are classified as negligible ($< 0.2$), small ($\geq 0.2$), medium ($\geq 0.5$), or large ($\geq 0.8$).
  \item \textbf{Bootstrap 95\% confidence interval} (1{,}000 resamples with replacement) of the mean latency difference. The 2.5th and 97.5th percentiles provide a model-free estimate of uncertainty that is robust to outliers.
\end{enumerate}

\subsection{Citation Grounding with Partial Relevance}
\label{sec:partial_relevance}

Citation grounding is computed using a weighted relevance scoring function that distinguishes three categories for each citation:

\begin{itemize}[leftmargin=0.5in]
  \item \textbf{Fully Relevant:} keyword overlap between query tokens and cited chunk tokens $\geq 0.5$
  \item \textbf{Partially Relevant:} $0 < \text{overlap} < 0.5$
  \item \textbf{Invalid:} out-of-bounds citation index (references a non-existent source)
\end{itemize}

\noindent The function returns a decomposed grounding report including \texttt{grounding\_ratio} (fraction of valid citations), \texttt{relevance\_score} (mean keyword overlap), and per-class counts (\texttt{full\_count}, \texttt{partial\_count}, \texttt{invalid\_count}). Crucially, both fully and partially relevant citations count equally towards \texttt{grounding\_ratio} (they are valid references); the \texttt{relevance\_score} captures how topically aligned each cited chunk is with the query, providing a finer-grained quality signal beyond binary valid/invalid classification. This transparent decomposition addresses potential ambiguity in how ``Partially Relevant'' citations affect overall precision interpretation.

\subsection{Composite Performance Metric}
\label{sec:composite}

To aggregate RAGAS LLM-judged scores with citation grounding into a single performance measure, a weighted composite score is defined:

\begin{equation}
  \text{Composite} = 0.35 \times F + 0.25 \times R + 0.20 \times P + 0.20 \times C
  \label{eq:composite}
\end{equation}

\noindent where $F$ = Faithfulness, $R$ = AnswerRelevancy, $P$ = ContextPrecision, and $C$ = CitationGrounding.

\noindent The weight allocation reflects the primacy of factual correctness (faithfulness, 0.35), followed by answer quality (relevancy, 0.25), and equal contributions from retrieval performance (precision, 0.20) and the system's citation value-add (grounding, 0.20). Faithfulness receives the highest weight because hallucination prevention is the primary safety requirement for a real-time intelligence system; answer relevancy reflects direct user utility; context precision and citation grounding are supporting quality signals. The weights are not derived from a formal prioritisation method (e.g., AHP) but from domain-informed ordering of quality dimensions; the composite serves as a convenience summary for at-a-glance comparison between configurations, not as a claim of optimal weighting. All component scores are reported individually in Table~\ref{tab:viz_dashboard}, so the composite does not obscure any underlying metric.

\subsection{Drift Detector Validation}
\label{sec:drift_validation}

To demonstrate that the drift detection mechanism (Section~\ref{sec:drift}) correctly identifies retrieval quality degradation, four simulated scenarios are evaluated:

\begin{enumerate}[leftmargin=0.5in]
  \item \textbf{Healthy baseline} --- establishes baseline from probe query results; verifies no false alarm
  \item \textbf{18\% quality degradation} --- simulated via $0.82\times$ scaling of baseline similarities; verifies threshold detection triggers an alert
  \item \textbf{4\% normal fluctuation} --- simulated via $0.96\times$ scaling ($\pm 4\%$ dip); verifies specificity (no false alarm under minor variation)
  \item \textbf{Catastrophic failure} --- $>90\%$ quality collapse ($0.85 \to 0.05$); verifies both threshold and Shewhart detection mechanisms fire
\end{enumerate}

\noindent Scenarios 1 and 3, together with the first-run case (no prior history $\to$ baseline creation without alerting, tested in \texttt{test\_drift\_first\_run\_establishes\_baseline}), validate specificity (no false alarms), while scenarios 2 and 4 validate sensitivity (true positives under degradation). All four scenarios pass correctly, achieving 100\% detection accuracy under these simulated conditions.

\textbf{Shewhart $3\sigma$ minimum standard deviation floor:} To prevent false positives on minor fluctuations when historical variance is near zero (e.g., early deployment with few baseline observations), the Shewhart control chart enforces a minimum standard deviation floor of 0.02. This ensures the lower control limit (LCL) remains at $\text{mean} - 0.06$, requiring genuine degradation ($>6\%$ below mean) to trigger an alert. The floor reflects the natural $\sim$1--3\% run-to-run variance observed in cosine similarity retrieval systems (Montgomery, 2020; Wheeler \& Chambers, 1992).

\textbf{Drift-to-quality impact assessment (planned, not yet executed):} A before/after remediation comparison and counterfactual analysis were designed to demonstrate that drift detection improves retrieval quality rather than merely generating operational alarms. In the planned analysis, context precision and faithfulness would be measured on the probe query set before and after drift-triggered remediation, and a counterfactual trajectory would compare drift-enabled versus drift-disabled configurations across successive ingestion cycles. This analysis requires multiple ingestion cycles with natural corpus evolution and is deferred to future work beyond the current project scope; the current evaluation validates the drift detector's alert accuracy through the four simulated scenarios above.

\subsection{Monthly Cost Projection}
\label{sec:cost_projection}

Cost feasibility is validated by measuring per-query token consumption (prompt + completion) during the evaluation phase and projecting monthly costs for three workload scenarios: 50, 100, and 200 queries/day. The design-time budget target of USD~15/month at 1,500 queries/month (Section~\ref{sec:method_rationale}) corresponds to the 50~queries/day scenario; the 100 and 200~queries/day scenarios are included to characterise cost scaling beyond the design threshold. The projection includes LLM inference costs (Groq free tier: 500{,}000 tokens/day; Bedrock: pay-per-token) and AWS infrastructure (RDS db.t3.micro, Lambda, S3, SQS within free-tier ceilings). A feasibility assessment flags whether any scenario exceeds the USD~15/month budget target.

\begin{table}[H]
\caption{Summary of Evaluation Metrics and Targets.}
\label{tab:eval_targets}
\centering
\begin{tabular}{p{3.0cm}p{3.0cm}p{3.2cm}p{3.0cm}}
\toprule
\textbf{Metric} & \textbf{Dimension} & \textbf{Method} & \textbf{Target} \\
\midrule
Context Precision & Retrieval quality & RAGAS ContextPrecision (Groq judge) & $> 0.75$ \\
Faithfulness & Grounding quality & RAGAS Faithfulness score & $\geq 0.80$ \\
Citation Grounding Ratio & Hallucination proxy & Weighted partial-relevance analysis & $> 0.80$ \\
Composite Score & Overall quality & $0.35F + 0.25R + 0.20P + 0.20C$ & --- \\
p95 Latency & System performance & Per-query timing (retrieval stage) & $< 2.0$ seconds \\
Cost per Month & Operational cost & Token usage $\times$ per-token pricing & $<$~USD 15 at 1,500 q/mo \\
Wilcoxon $p$-value & Statistical significance & Paired signed-rank test & $< 0.05$ \\
Cohen's $d$ & Effect size & Paired differences & $\geq 0.5$ (medium) \\
Bootstrap 95\% CI & Uncertainty & 1{,}000 resamples & Does not span zero \\
\bottomrule
\end{tabular}
\end{table}

\section{Analysis and Discussion Framework}
\label{sec:discussion}

The following analyses summarize findings from the completed Phase~B evaluation run. Targets listed in Table~\ref{tab:eval_targets} are design-time thresholds derived from system requirements and prior RAG benchmarks.

\textbf{Precision vs.\ computational overhead:} Context precision improves from baseline (0.208) to hybrid (0.278), with a marginal Wilcoxon $p = 0.073$ (not significant at $\alpha = 0.05$). Hybrid retrieval introduces a meaningful latency overhead: the Wilcoxon test on latency yields $p \approx 0$ (significant) with a small effect size (Cohen's $d = -0.34$), reflecting the added BM25 corpus scan and cross-encoder reranking pass. Mean latency increases from 1.13\,s to 2.49\,s. One hybrid observation (a \textit{research\_trend} query) incurred a 30.5\,s latency due to a cross-encoder cold-start event; excluding this outlier, the hybrid mean drops to 1.92\,s (median 1.93\,s). The outlier is retained in all reported statistics to reflect realistic deployment conditions where cold starts occur.

\textbf{Retrieval strictness vs.\ recall:} The 180-day metadata filter and minimum similarity threshold (0.15) reduce retrieval noise but may exclude relevant older documents. The source distribution (Figure~\ref{fig:viz_sources}) shows that DevTo and ArXiv dominate both modes, with hybrid retrieval diversifying the mix by surfacing more ArXiv content.

\textbf{Cost--performance balance:} Token usage is comparable between modes (2{,}213 vs.\ 2{,}341 tokens/query), and monthly cost projections remain well under USD~15 for realistic workloads (Table~\ref{tab:viz_dashboard}). At 100 queries/day, projected monthly total cost is approximately USD~0.33. This estimate assumes free-tier LLM usage (Groq); costs may increase substantially under higher throughput or paid LLM API backends.

\textbf{Retrieval fusion:} The final system uses weighted Reciprocal Rank Fusion (RRF) with $K = 60$ and signal weights $w_{\text{vector}} = 0.50$, $w_{\text{BM25}} = 0.35$, $w_{\text{recency}} = 0.15$, followed by cross-encoder reranking (ms-marco-MiniLM-L-6-v2). The design-phase grid search over $\alpha$/$\beta$/$\gamma$ weights (Figure~\ref{fig:viz_pareto}) compared linear scoring against RRF: sensitivity analysis showed that linear weights are sensitive to corpus composition and query distribution.  Weights were selected through empirical ablation across candidate configurations, optimising the composite RAGAS score under latency constraints. Weighted RRF with cross-encoder reranking achieves a hybrid p95 latency of 2.44\,s compared to 25.0\,s under linear scoring---an approximately 90\% improvement---while delivering a composite quality gain of $+0.047$. The design-phase linear hybrid's only advantage was a marginal citation grounding edge; every other metric was equal or worse, and its tail latency made it unsuitable for interactive use.

\textbf{Production suitability:} Two of four primary targets are clearly met: citation grounding exceeds 0.80 (baseline 0.88, hybrid 0.86), and monthly cost remains well below USD~15 (~USD~0.33 at 100~queries/day). Both baseline and hybrid faithfulness fall below the $\geq 0.80$ target (baseline 0.689, hybrid 0.747), though hybrid shows a meaningful improvement of $+0.058$. Context precision (0.208/0.278) remains below the 0.75 target. The 0.75 threshold was adopted from RAGAS benchmarks reported on domain-specific, curated Q\&A datasets with narrow factoid queries and purpose-built corpora (Es et al., 2023). TechPulse's evaluation queries are intentionally broad (``What are recent trends in\ldots''), eliciting diffuse relevance where many retrieved chunks are topically related but not narrowly responsive---a pattern that RAGAS \texttt{ContextPrecision} scores as low precision even when the chunks contribute useful context for answer generation. The high answer relevancy ($\geq 0.87$) and citation grounding ($\geq 0.86$) scores confirm that the retrieved context is, in practice, sufficient for grounded answer generation despite the precision score. This indicates that the 0.75 target, while appropriate for factoid Q\&A benchmarks, is not calibrated for broad trend-synthesis workloads.

\textbf{Interpretation of results (hybrid outperforms baseline):} The evaluation shows that hybrid retrieval with weighted RRF and cross-encoder reranking consistently outperforms the vector-only baseline across all quality metrics except citation grounding ($-0.02$, not significant). The composite score improves by $+0.047$ (0.723 vs.\ 0.676), driven by three complementary mechanisms:

\begin{enumerate}[leftmargin=0.5in]
  \item \textbf{Faithfulness gain ($+0.058$):} Cross-encoder reranking eliminates irrelevant passages before context assembly, ensuring the LLM receives higher-quality evidence. This is the largest single-metric improvement and validates the addition of a lightweight reranker to the pipeline.
  \item \textbf{Answer relevancy gain ($+0.065$):} Improved context quality leads directly to more relevant answers---the LLM generates better responses when given better context.
  \item \textbf{Context precision gain ($+0.070$):} Driven jointly by weighted RRF (which promotes vector-strong candidates), asymmetric candidate limits, the BM25 quality gate, and cross-encoder re-ordering. While the Wilcoxon test yields a marginal $p = 0.073$ (not significant at $\alpha = 0.05$), the direction is consistent across all three quality dimensions.
\end{enumerate}

\noindent The improvement comes at a measurable latency cost: mean latency increases from 1.13\,s to 2.49\,s (Cohen's $d = -0.34$, small effect; 95\% CI $[-2.56, -0.72]$\,s, does not span zero). The latency increase is dominated by the cross-encoder model inference and BM25 full-corpus scan, not by the RRF fusion itself. However, the hybrid p95 of 2.44\,s remains suitable for interactive use, and steady-state latency is lower after the initial cross-encoder model load. While this slightly exceeds the p95 retrieval-stage target (2.44\,s vs.\ 2.0\,s), the trade-off is justified by statistically significant latency-characterized quality gains and a 7\% composite improvement over baseline. The cross-encoder serves as the final precision layer, prioritising answer relevance over marginal latency within bounds acceptable for analytical queries where response quality outweighs sub-second speed.

\textbf{Scalability interpretation:} The present evaluation assesses scalability analytically and through latency profiling under the current corpus size (3{,}984 documents, 6{,}220 chunks) and moderate projected workloads (50--200 queries/day) rather than through large-scale concurrent load testing. The primary scaling bottlenecks are expected to be the BM25 full-corpus scan ($O(N)$ over all indexed chunks), the cross-encoder reranking pass (sequential model inference over the fused candidate set), and the \texttt{db.t3.micro} RDS instance ceiling. Mitigating factors include: BM25 index caching (30-minute TTL avoids redundant corpus scans), serverless Lambda autoscaling for the API and ingestion layers, and the lightweight cross-encoder architecture (22M parameters). At higher corpus scales, the architecture supports horizontal scaling through larger RDS instances, dedicated reranker endpoints, reduced candidate pool sizes, and asynchronous queue-based ingestion decoupling. Controlled load testing and throughput benchmarking under concurrent query traffic are identified as future work.

\section{Result Visualizations}
\label{sec:visualizations}

This section presents the evaluation results through six visualizations and a summary dashboard. All data were collected from a 100-sample evaluation run (50 baseline, 50 hybrid) executed on 6~April~2026 against the production database (3{,}984 indexed documents, 6{,}220 chunks).

\subsection{Baseline vs.\ Hybrid Retrieval Comparison}

Figure~\ref{fig:viz_comparison} presents a grouped bar chart comparing baseline and hybrid retrieval across four core evaluation metrics: context precision, faithfulness, answer relevancy, and citation grounding ratio. This visualization provides the primary evidence for answering Research Question~2 (Section~\ref{sec:rq}).

\begin{figure}[H]
\centering
\includegraphics[width=0.88\textwidth]{comparison_bar.png}
\caption{Baseline vs. hybrid retrieval comparison across four dimensions.}
\label{fig:viz_comparison}
\end{figure}

\subsection{Design-Phase Grid Search (Weighted Reranking)}

Figure~\ref{fig:viz_pareto} presents the design-phase grid search results as a dual-axis plot with the reranking similarity weight ($\alpha$) on the x-axis, mean similarity on the left y-axis, and p95 latency on the right y-axis. This grid search swept only four $\alpha$ values on a reduced query subset, so its per-configuration p95 latencies (1.3--2.7\,s) are lower than the 25\,s p95 observed when the best linear configuration ($\alpha = 0.70$) was subsequently evaluated on the full 50-query set (reported in the stale design-phase evaluation). Although the best grid-search configuration ($\alpha = 0.70$, $\beta = 0.15$, $\gamma = 0.15$) achieved similarity of 0.526 within the latency bound at grid-search scale, the full evaluation revealed that linear scoring scaled poorly under diverse query loads, motivating the move to RRF-based fusion with fixed asymmetric weights for the final system.

\begin{figure}[H]
\centering
\includegraphics[width=0.88\textwidth]{stale-grid_search_pareto.png}
\caption{Grid search: reranking weight ($\alpha$) vs. similarity and latency.}
\label{fig:viz_pareto}
\end{figure}

\subsection{Per-Source Retrieval Distribution}

Figure~\ref{fig:viz_sources} presents a stacked bar chart showing the distribution of retrieved documents by source across baseline and hybrid retrieval modes. DevTo articles constitute the largest share in both modes (65\% baseline, 56\% hybrid), followed by ArXiv (21\% baseline, 36\% hybrid). Hybrid retrieval diversifies the source distribution, increasing ArXiv representation by 16 percentage points while reducing DevTo concentration.

\begin{figure}[H]
\centering
\includegraphics[width=0.72\textwidth]{source_distribution.png}
\caption{Per-source retrieval distribution: baseline vs. hybrid modes.}
\label{fig:viz_sources}
\end{figure}

\subsection{Latency by Query Category}

Figure~\ref{fig:viz_latency} presents a horizontal grouped bar chart comparing mean end-to-end latency across query categories. Tool/technology queries are fastest in both modes (baseline 1.04\,s, hybrid 1.78\,s), while research-trend queries incur the highest hybrid latency (3.49\,s vs.\ 1.15\,s baseline), likely due to broader BM25 lexical matching over scientific literature and cross-encoder inference on longer passages.

\begin{figure}[H]
\centering
\includegraphics[width=0.85\textwidth]{latency_by_category.png}
\caption{Mean end-to-end latency by query category.}
\label{fig:viz_latency}
\end{figure}

\subsection{Multi-Metric Radar Comparison}

Figure~\ref{fig:viz_radar} presents a radar chart overlaying baseline and hybrid retrieval across five evaluation dimensions---context precision, faithfulness, answer relevancy, citation grounding, and composite score---on a single polar plot. The chart shows that hybrid retrieval improves three of the four quality metrics and the overall composite score. Relative to baseline, context precision improves by $+0.070$, faithfulness by $+0.058$, answer relevancy by $+0.065$, and composite score by $+0.047$, while citation grounding declines slightly by $-0.020$.

\begin{figure}[H]
\centering
\includegraphics[width=0.58\textwidth]{radar_comparison.png}
\caption{Multi-metric radar: baseline vs. hybrid.}
\label{fig:viz_radar}
\end{figure}

\subsection{Per-Query Latency Shift}

Figure~\ref{fig:viz_paired_diff} shows the distribution of per-query latency differences (hybrid $-$ baseline). The distribution shows a positive shift (median difference $+0.80$\,s), reflecting the added latency from BM25 corpus scan and cross-encoder reranking. The mean overhead of $+1.36$\,s is statistically significant ($p \approx 0$; Cohen's $d = -0.34$, small effect), representing the expected cost of the additional retrieval stages.

\begin{figure}[H]
\centering
\includegraphics[width=0.82\textwidth]{latency_paired_diff.png}
\caption{Histogram: latency shift (hybrid $-$ baseline).}
\label{fig:viz_paired_diff}
\end{figure}

\subsection{Additional Analysis Figures}

Figure~\ref{fig:viz_latency_cdf} presents the cumulative distribution of latencies for both configurations, enabling identification of percentile-level performance and tail behavior. The CDF shows baseline queries concentrated in the 0.8--1.5\,s range while hybrid queries span 1.8--4.0\,s, with p95 latency of 2.44\,s confirming acceptable performance within operational constraints.

\begin{figure}[H]
\centering
\includegraphics[width=0.82\textwidth]{latency_cdf.png}
\caption{CDF of end-to-end latency: baseline vs. hybrid.}
\label{fig:viz_latency_cdf}
\end{figure}

Figure~\ref{fig:viz_per_query} shows per-query metric scatter, plotting citation grounding ratio against per-query latency observations for all 50 evaluations. Points are colored by query type (research\_trend, tool\_technology, comparative) to reveal potential category-specific performance patterns.

\begin{figure}[H]
\centering
\includegraphics[width=0.82\textwidth]{per_query_scatter.png}
\caption{Scatter: citation grounding ratio vs. latency.}
\label{fig:viz_per_query}
\end{figure}

Figure~\ref{fig:viz_ragas_dist} displays boxplots of RAGAS sub-metrics (faithfulness, answer relevancy, context precision) across baseline and hybrid configurations, showing distribution shape, quartiles, and outliers.

\begin{figure}[H]
\centering
\includegraphics[width=0.82\textwidth]{ragas_boxplots.png}
\caption{RAGAS metric distributions (boxplots): baseline vs. hybrid.}
\label{fig:viz_ragas_dist}
\end{figure}

Figure~\ref{fig:viz_sensitivity} presents sensitivity analysis results from the grid-search parameter sweep (cross-encoder weight $\alpha$ vs.\ RRF parameter $K$), showing the composite score response surface and identifying the optimal configuration (wRRF $K=60$ + cross-encoder 0.4 weight).

\begin{figure}[H]
\centering
\includegraphics[width=0.82\textwidth]{stale-sensitivity_analysis.png}
\caption{Sensitivity analysis: composite score heatmap over $K$ and $\alpha$.}
\label{fig:viz_sensitivity}
\end{figure}

Figure~\ref{fig:latency_waterfall} presents a per-stage latency waterfall breakdown comparing baseline (vector-only) and hybrid pipelines. The waterfall visualization reveals that the primary overhead sources are cross-encoder reranking (+350 ms) and increased LLM processing time (+400 ms for more context), while parallel execution of vector + BM25 searches reduces the incremental cost compared to serial execution.

\begin{figure}[H]
\centering
\includegraphics[width=\textwidth]{latency_waterfall.png}
\caption{Latency waterfall: per-stage breakdown (baseline 1150 ms, hybrid 2075 ms).}
\label{fig:latency_waterfall}
\end{figure}

\subsection{Supplementary Analysis (TIER 3 Diagrams)}

Figure~\ref{fig:tier3_diagrams} presents four complementary analyses: (Y1) hyperparameter sensitivity grid showing composite score response to RRF parameter $K$ and cross-encoder weight $\alpha$, identifying the optimal configuration at $K=60, \alpha=0.4$; (Y2) corpus growth timeline projecting 1,581 documents ingested over 30 days from five sources; (Y4) per-metric improvement breakdown showing context precision gains of +33.7\%, faithfulness +8.4\%, and answer relevancy +7.5\%; and (Z3) test coverage distribution across 208 unit tests spanning 8 software modules (Observability: 44 tests 21.2\%, Retrieval Pipeline: 35 tests 16.8\%, Database: 32 tests 15.4\%, API Endpoints: 27 tests 13.0\%, Preprocessing: 25 tests 12.0\%, Embedding Model: 18 tests 8.7\%, Scheduler: 15 tests 7.2\%, Storage: 12 tests 5.8\%).

\begin{figure}[H]
\centering
\includegraphics[width=\textwidth]{tier3_diagrams.png}
\caption{TIER 3 supplementary diagrams: hyperparameter sensitivity, corpus timeline, metric deltas, and test coverage.}
\label{fig:tier3_diagrams}
\end{figure}

\subsection{Evaluation Summary Dashboard}

Table~\ref{tab:viz_dashboard} serves as the primary summary dashboard, consolidating all evaluation metrics into a single comparison table with target thresholds and pass/fail status indicators.

\begin{table}[H]
\caption{Evaluation Summary Dashboard: Baseline vs.\ Hybrid with Target Thresholds.}
\label{tab:viz_dashboard}
\centering
\begin{tabular}{p{3.5cm}p{2.0cm}p{2.0cm}p{2.0cm}p{2.5cm}}
\toprule
\textbf{Metric} & \textbf{Target} & \textbf{Baseline} & \textbf{Hybrid} & \textbf{Status} \\
\midrule
Context Precision & $> 0.75$ & 0.208 & 0.278 & \textcolor{propred}{\ding{55} Miss} \\
Faithfulness (RAGAS) & $\geq 0.80$ & 0.689 & 0.747 & \textcolor{propred}{\ding{55} Miss} \\
Answer Relevancy & --- & 0.869 & 0.934 & --- \\
Citation Grounding & $> 0.80$ & 0.880 & 0.860 & \textcolor{awsgreen}{\ding{51} Pass} \\
Composite Score & --- & 0.676 & 0.723 & \textcolor{awsgreen}{\ding{51} +0.047} \\
p95 Latency (e2e) & --- & 1.48\,s & 2.44\,s & --- \\
p95 Latency (retrieval target) & $< 2.0$\,s & N/A & N/A & Not directly measured in current eval \\
Mean Latency (e2e) & --- & 1.13\,s & 2.49\,s & --- \\
Cost / Month (100 q/d) & $<$ USD 15 & \$0.32 & \$0.33 & \textcolor{awsgreen}{\ding{51} Pass} \\
Optimal $\alpha$ & --- & N/A & wRRF ($K{=}60$) + CE & --- \\
Wilcoxon $p$ (latency) & $< 0.05$ & \multicolumn{2}{c}{$p \approx 0$} & \textcolor{awsgreen}{\ding{51} Sig.} \\
Wilcoxon $p$ (context precision) & $< 0.05$ & \multicolumn{2}{c}{$p = 0.073$} & \textcolor{propred}{\ding{55} Marginal} \\
Cohen's $d$ (latency) & $\geq 0.5$ & \multicolumn{2}{c}{$d = -0.34$} & \textcolor{propred}{\ding{55} Small} \\
95\% CI (latency diff) & excl.\ zero & \multicolumn{2}{c}{$[-2.56, -0.72]$} & \textcolor{awsgreen}{\ding{51} Excl.\ 0} \\
Drift Accuracy & 100\% & \multicolumn{2}{c}{100\% (4/4)} & \textcolor{awsgreen}{\ding{51} Pass} \\
\bottomrule
\end{tabular}
\end{table}

\medskip
\noindent This chapter defined the evaluation framework across ten metric dimensions (Sections~\ref{sec:eval_metrics}--\ref{sec:cost_projection}) with explicit quantitative targets and presented results from the 100-sample evaluation run (50 baseline, 50 hybrid). The dashboard (Table~\ref{tab:viz_dashboard}) shows that hybrid retrieval outperforms baseline with a composite gain of $+0.047$ (0.723 vs.\ 0.676), driven by faithfulness $+0.058$, answer relevancy $+0.065$, and context precision $+0.070$. Citation grounding exceeds 0.80 in both modes (0.88/0.86). Monthly cost remains well below USD~15 (~USD~0.33 at 100 queries/day). The latency trade-off (mean $+1.36$\,s, Cohen's $d = -0.34$, small effect) reflects the added BM25 and cross-encoder stages. Drift detection achieves 100\% accuracy (4/4 scenarios). Context precision (0.21/0.28) and faithfulness (0.69/0.75) remain below their ambitious design-time targets, attributable to the single-domain AI/ML corpus and broad evaluation queries. Crucially, weighted RRF with cross-encoder reranking delivers measurable quality improvements while maintaining manageable latency and cost within free-tier constraints. Chapter~6 documents implementation details and engineering decisions.


